% TOP_PROBLEM: {"id": 2741, "title": "Kirby Problem 1.82", "category_id": 7, "category": {"id": 7, "name": "topology", "display_name": "Topology", "description": "Properties preserved under continuous deformations.", "slug": "topology", "order_index": 7, "created_at": "2026-07-31T15:26:25.671Z"}, "status": "open", "problem_number": "KP-1.82", "set_id": 11}
% FIRST_POSTED: 2026-10-01
% ENTRY_KIND: statement_only
% UnsolvedMath ID 2741; source catalogue code KP-1.82
% !TeX program = lualatex
% Definition and sources only; this file is not an LLM solution attempt.
\documentclass[11pt,a4paper]{article}
\usepackage[margin=25mm]{geometry}
\usepackage{fontspec}
\setmainfont{lmroman10-regular.otf}[BoldFont=lmroman10-bold.otf,
 ItalicFont=lmroman10-italic.otf,BoldItalicFont=lmroman10-bolditalic.otf]
\usepackage{amsmath,amssymb,mathtools}
\usepackage{unicode-math}
\setmathfont{latinmodern-math.otf}
\AtBeginDocument{\let\setminus\smallsetminus}
\usepackage{xurl,hyperref}
\hypersetup{colorlinks=true,urlcolor=blue}
\setcounter{secnumdepth}{0}
\setlength{\parindent}{0pt}
\setlength{\parskip}{5pt}
\setlength{\emergencystretch}{3em}
\begin{document}
\section*{Kirby Problem 1.82}\label{2741}
{\small UnsolvedMath ID 2741 · KP-1.82 · Topology · Kirby's Problems in Low-Dimensional Topology}
\subsection{Definitions and mathematical statement}
An oriented knot or link is positive if it admits a regular diagram in which
each crossing has positive oriented sign under the usual right-hand convention.
This is an existence condition on diagrams: another diagram of the same link
may contain negative crossings. The exterior of a link $L$ is
$X_L=S^3\setminus\operatorname{int}\nu L$, together with its meridians and
longitudes as peripheral data.

The first question seeks an intrinsic characterization of positive knots.
Give necessary and sufficient conditions phrased in the topology or geometry
of the knot and its exterior, such as suitable spanning surfaces or group and
peripheral structure, that do not merely repeat the existence of a positive
diagram.

The second question asks for a simple explicit calculus relating positive
diagrams of the same oriented knot or link. Specify moves that preserve the
link type and positivity, and prove that any two positive diagrams of that
type are connected by finitely many such moves and planar isotopies. A complete
calculus must allow positive diagrams with different numbers of crossings;
the problem does not assume that every positive diagram is crossing-minimal.
\subsection{Short English statement}
Characterize positive knots without choosing a positive diagram, and find moves connecting all positive diagrams of a given link.
\subsection{Sources}
\begin{itemize}
\item[S1] ULAM.AI and the UnsolvedMath Contributors, supplied problems.json, record 2741 (KP-1.82), Kirby's Problems in Low-Dimensional Topology. Catalogue identity and source transcription; CC BY 4.0.
\url{https://huggingface.co/datasets/ulamai/UnsolvedMath}.
\item[S2] K3: A New Problem List in Low-Dimensional Topology, Problem 1.82. Expanded formulation of the supplied catalogue statement.
\url{https://math.berkeley.edu/sites/default/files/surv-295-ruberman-watermarked-author-pdf.pdf}.
\end{itemize}
\end{document}
